\textbf{\color{RoyalBlue}Erkennen}
= woran man den Aufgabentyp erkennt, \textbf{\color{RedViolet}Achtung}
= häufige Fehler.


\section{Formel- und Methodenübersicht zu den Präsenzaufgaben (erstellt von Claude, NOCH NICHT FERTIG ÜBERPRÜFT)}\label{sec:formel--und-methodenubersicht-zu-den-prasenzaufgaben-(erstellt-von-claude-noch-nicht-fertig-uberpruft)}

% ==================================================================

\subsection{Zuordnung: Aufgabe $\to$ Aufgabentyp}\label{subsec:zuordnung:-aufgabe-$to$-aufgabentyp}
% ==================================================================

\begin{mytable}
    \begin{tabular}{p{2.3cm} p{6.2cm} p{6.4cm}}
        \hline
        \textbf{Aufgabe} & \textbf{Thema / Werkzeug}                 & \textbf{Kernbegriffe}                                \\
        \hline
        1.0.1            & Stetige Gleichverteilung, Ziffern         & Lebesgue-Maß, Längen, $\pe(\text{Punkt})=0$          \\
        1.0.2, 1.1.6     & Modellierung, abzählbar unendl.\ $\eraum$ & Geometrische Struktur, Summe der Fortschritte        \\
        1.0.3            & Kombinatorik / Permutationen              & Laplace, Fixpunkte, Inklusion-Exklusion, Derangement \\
        1.1.4, 1.1.5     & Münzwurf, Laplace-Produktraum             & $\{K,Z\}^n$, Abzählen, ``mind.\ einmal''             \\
        1.2.4            & Dichten (stetig)                          & Normierung, $\pe$ von Intervallen                    \\
        1.3.4            & Zufallsvariablen \& Verteilung            & Verteilung, Verteilungsfkt., gemeinsame Vert.      \\
        1.4.4            & Unabh.\ von Ereignissen                   & Produktraum, paarweise vs.\ gemeinsame Unabh.      \\
        1.4.5            & 2D-Gleichverteilung                       & Randverteilung, Unabh.\ von ZV                       \\
        1.5.4            & Gemeinsame Dichte                         & Doppelintegral, Randdichten, Faktorisierung          \\
        1.5.5            & Unabh.\ (Theorie), $\max$ von ZV          & Funktionen unabh.\ ZV                                \\
        1.6.4, 1.6.6     & Urnen, Bayes                              & Totale W., Satz von Bayes, ohne Zurücklegen          \\
        1.6.5            & Bedingte W.\ (Theorie)                    & Rechenregeln bedingter W.                          \\
        2.1.4, 2.1.5     & $\E{}$, $\Var$ diskret                    & Vert.-tabelle, lineare Transformation $aX+b$         \\
        2.1.6            & $\E{}$, $\Var$ stetig                     & Momente per Integral                                 \\
        2.2.4, 2.3.6     & Binomial + Normalapprox.
                       & de Moivre--Laplace, Stetigkeitskorrektur             \\
        2.2.5            & Binomial, Tschebyscheff                   & Konzentration, Gesetz großer Zahlen                  \\
        2.2.6            & Geometrische Verteilung                   & Wartezeit, $\E{},\Var$, Tschebyscheff                \\
        2.3.4, 2.3.5     & Poisson-Approximation                     & seltene Ereignisse, $\lambda=np$                     \\
        2.4.4            & Schätztheorie                             & Erwartungstreue, Risiko (MSE), Konsistenz            \\
        2.4.5            & Deskriptive Statistik                     & Mittel, Median, Modus, Streuung                      \\
        2.5.4, 2.5.5     & Konfidenzintervalle                       & Anteil $p$, Wald-Intervall                           \\
        2.6.4, 2.6.5     & Hypothesentests                           & einseitiger Anteilstest, Fehler 1./2.\ Art           \\
        \hline
    \end{tabular}
\end{mytable}

% ==================================================================

\subsection{Wahrscheinlichkeitsräume modellieren}\label{subsec:wahrscheinlichkeitsraume-modellieren}
% ==================================================================
Fast jede Aufgabe beginnt mit ``Geben Sie einen geeigneten Wahrscheinlichkeitsraum
$(\eraum,\pe)$ an''.
Drei Grundtypen deckt der gesamte Kurs ab.

\subsubsection{Typ A -- Endlicher Laplace-Raum (Gleichverteilung)}
Alle Ergebnisse gleich wahrscheinlich.Verwenden bei Münze, Würfel,
Permutationen, Ziehen mit/ohne Zurücklegen.
\[
    \eraum \text{ endlich},\qquad
    \pe(A)=\frac{\#A}{\#\eraum}=\frac{\text{günstige}}{\text{mögliche}}.
\]
\faustregel{``fair'', ``rein zufällig'', symmetrisch, endlich viele Ausgänge
    $\Rightarrow$ zählen.}
\warnung{Der Laplace-Ansatz gilt \emph{nur}, wenn die Ergebnisse in $\eraum$
    gleichwahrscheinlich sind.Beispiel Augensumme dreier Würfel: $\eraum=\{1,\dots,6\}^3$
    ist Laplace, aber die \emph{Summe} $Z$ ist es \textbf{nicht} -- also im Grundraum
    zählen, nicht auf der Summenebene.}

\subsubsection{Typ B -- Diskreter, abzählbar unendlicher Raum}
Wartezeiten (Aufg.\ 1.0.2, 1.1.6, 2.2.6): ``wiederhole, bis \dots''.
\[
    \eraum=\nz \text{ (oder } \nz_0\text{)},\qquad
    \pe(\{k\})=p_k,\quad \sum_{k}p_k=1.
\]
\warnung{Bei M\"ensch-ärgere-dich-nicht: die einmalige Würfe bilden einen
endlichen Raum, aber \emph{``bis zum ersten Nicht-6''} erzeugt einen unendlichen
Raum.Man kann denselben Sachverhalt oft auf mehreren $\eraum$ modellieren
    (Frage ``weitere sinnvolle Wahlen''): fein (ganze Würfelfolge) oder grob
    (nur die Endsumme).Wichtig ist, dass die interessierende Größe messbar bleibt.}

\subsubsection{Typ C -- Überabzählbarer Raum (kontinuierlich)}
Zahl aus $[0,1]$, Ankunftszeiten (Aufg.\ 1.0.1, 1.4.5, 2.3.6).
\[
    \eraum\subseteq\rz^d,\quad \pe(A)=\int_A f(x)dx
    \quad(\text{Gleichverteilung: } f=\frac{1}{Vol(\eraum)}\mathbf{1}_{\eraum}).
\]
\begin{reg}
    \textbf{Erkennen:} kontinuierlicher Wertebereich, ``uniform / rein
    zufällig aus einem Intervall/Gebiet''.Wahrscheinlichkeit = \textbf{Länge /
    Fläche / Volumen} des Ereignisses geteilt durch die des Grundgebiets.
\end{reg}
\warnung{$\pe(\{x\})=0$ für jeden einzelnen Punkt.``Die Zahl ist \emph{exakt}
    $\frac17$'' hat also Wahrscheinlichkeit $0$; ``höchstens $\frac1{100}$ entfernt''
    ist ein Intervall der Länge $\frac{2}{100}$.}

% ==================================================================

\subsection{Kombinatorik und Permutationen}\label{subsec:kombinatorik-und-permutationen}
% ==================================================================
Grundlage für alle Laplace-Aufgaben (1.0.3, 1.1, 1.4.4).

\begin{mytable}
    \begin{tabular}{@{}l l l@{}}
        \hline
        & \textbf{mit Zurücklegen} & \textbf{ohne Zurücklegen} \\
        \hline
        \textbf{mit Reihenfolge}  & $n^k$                    & $\dfrac{n!}{(n-k)!}$      \\
        \textbf{ohne Reihenfolge} & $\dbinom{n+k-1}{k}$      & $\dbinom{n}{k}$           \\
        \hline
    \end{tabular}
\end{mytable}

\noindent Permutationen: $\#\{\text{Bijektionen von }\{1,\dots,n\}\}=n!$.

\paragraph{Fixpunkte (Aufg.\ 1.0.3).} Mit $A_l=\{\omega:\omega(l)=l\}$ und
Inklusion-Exklusion:
\[
    \pe(\text{genau } k \text{ Fixpunkte})
    =\frac{1}{k!}\sum_{j=0}^{n-k}\frac{(-1)^j}{j!}
    \;\xrightarrow[n\to\infty]{}\; \frac{e^{-1}}{k!}.
\]
Speziell ``kein Fixpunkt'' (Derangement): $D_n=n!\sum_{j=0}^n\frac{(-1)^j}{j!}$.
\warnung{``mindestens $k$'' ($C_k$) $\neq$ ``genau $k$'' ($B_k$).Zusammenhang
    $\pe(B_k)=\pe(C_k)-\pe(C_{k+1})$.Bei Inklusion-Exklusion die
    \emph{Schnitte} $\pe(A_{l_1}\cap\dots\cap A_{l_j})=\frac{(n-j)!}{n!}$ verwenden,
    nicht na\"iv multiplizieren.}

% ==================================================================

\subsection{Dichten stetiger Verteilungen}\label{subsec:dichten-stetiger-verteilungen}
% ==================================================================
Aufg.\ 1.2.4, 1.5.4, 2.1.6.

\paragraph{Dichte-Eigenschaften.} $f\geq 0$ und
$\displaystyle\int_{-\infty}^{\infty} f(x)dx = 1$.
Normierungskonstante $c$ so bestimmen, dass das Integral $=1$ wird:
\[
    \int_{\text{Träger}} c\,g(x)dx = 1
    \;\Longrightarrow\;
    c=\Bigl(\int g\Bigr)^{-1}.
\]

\paragraph{Wahrscheinlichkeiten \& Verteilungsfunktion.}
\[
    \pe\bigl((a,b]\bigr)=\int_a^b f,\qquad
    F(x)=\pe(X\leq x)=\int_{-\infty}^x f(t)dt,\qquad f=F'.
\]
\warnung{$\pe(\{a\})=0$ für stetige Verteilungen: offene/geschlossene Grenzen
egal.Integrationsgrenzen immer auf den \textbf{Träger} $\{f>0\}$ einschränken --
z.B.\ liefert $\pe([1,\infty))$ bei Träger $[0,1]$ nur den Punkt $\{1\}$, also $0$.}

% ==================================================================

\subsection{Zufallsvariablen: Verteilung, Verteilungsfunktion, gemeinsame Verteilung}\label{subsec:zufallsvariablen:-verteilung-verteilungsfunktion-gemeinsame-verteilung}
% ==================================================================
Aufg.\ 1.3.4.

\begin{itemize}
    \item \textbf{Verteilung} (diskret): Tabelle $x\mapsto\pe(X=x)$ (Summe $=1$).
    \item \textbf{Verteilungsfunktion}: $F_X(t)=\pe(X\leq t)$; diskret eine
    rechtsstetige Treppenfunktion mit Sprüngen der Höhe $\pe(X=x)$.
    \item \textbf{Gemeinsame Verteilung}: $\pe(X=x,\,Y=y)$ als Matrix.
    \item \textbf{Randverteilung}: Zeilen-/Spaltensummen,
    $\pe(X=x)=\sum_y \pe(X=x,Y=y)$.
\end{itemize}
\warnung{Werte\emph{bereich} (Definitionsmenge der Argumente) $\neq$
    Bild\emph{bereich} $X(\eraum)$ (angenommene Werte).Bei ``mit Zurücklegen'' sind
    $X_1,X_2$ unabh.\ und die gemeinsame Verteilung ist das Produkt der Rand\-verteilungen;
    bei ``ohne Zurücklegen'' \emph{nicht}.$Z=\min\{X_1,X_2\}$ hängt von beiden ab,
    ist also i.\,A.\ \emph{nicht} unabh.\ von $X_1$.}

% ==================================================================

\subsection{Unabhängigkeit}\label{subsec:unabhangigkeit}
% ==================================================================

\subsubsection{Ereignisse (Aufg.\ 1.4.4)}
\[
    A,B \text{ unabh.} \iff \pe(A\cap B)=\pe(A)\,\pe(B).
\]
Gemeinsame (vollständige) Unabhängigkeit von $A_1,\dots,A_m$: die
Produktregel muss für \emph{jede} Teilmenge der Indizes gelten.
\warnung{\textbf{Paarweise Unabhängigkeit $\not\Rightarrow$ gemeinsame
Unabhängigkeit.} Bei ``$A,B,C$ unabh.?'' nach dem Prüfen der Paare zusätzlich
    $\pe(A\cap B\cap C)=\pe(A)\pe(B)\pe(C)$ nachrechnen.}

\subsubsection{Zufallsvariablen (Aufg.\ 1.4.5, 1.5.4, 1.5.5)}
\[
    X,Y \text{ unabh.} \iff F_{X,Y}(x,y)=F_X(x)F_Y(y)
    \iff f_{X,Y}(x,y)=f_X(x)f_Y(y).
\]
\begin{reg}
    \textbf{Erkennen (Dichte):} $X,Y$ unabh.\ $\iff$ die gemeinsame
    Dichte \emph{faktorisiert} \textbf{und} der Träger ein Rechteck (Produktmenge) ist.
    Ein ``schräger'' Träger wie $\{0\leq x\leq y\leq 1\}$ koppelt die Variablen
    $\Rightarrow$ \textbf{nicht} unabhängig, selbst wenn $f$ als Produkt aussieht.
\end{reg}

\paragraph{Funktionen unabh.\ ZV (1.5.5).} Sind $X,Y,Z$ unabh., so sind auch
$\max\{X,Y\}$ und $Z$ unabh.\ (disjunkte Variablengruppen).
\warnung{$\max\{X,Y\}$ und $\max\{X,Z\}$ sind i.\,A.\ \emph{nicht} unabh.\ --
sie teilen sich $X$.Gemeinsame Variablen zerstören Unabhängigkeit.}

% ==================================================================

\subsection{Bedingte Wahrscheinlichkeiten und Bayes}\label{subsec:bedingte-wahrscheinlichkeiten-und-bayes}
% ==================================================================
Aufg.\ 1.6.4--1.6.6.
\[
    \pe(A\mid B)=\frac{\pe(A\cap B)}{\pe(B)},\quad \pe(B)>0.
\]

\paragraph{Satz von der totalen Wahrscheinlichkeit} (Zerlegung $B_1,\dots,B_n$):
\[
    \pe(A)=\sum_i \pe(A\mid B_i)\,\pe(B_i).
\]

\paragraph{Satz von Bayes}
\[
    \pe(B_k\mid A)=\frac{\pe(A\mid B_k)\,\pe(B_k)}{\sum_i \pe(A\mid B_i)\,\pe(B_i)}.
\]
\begin{reg}
    \textbf{Erkennen:} Zweistufiges Experiment (erst Urne/Münze
    wählen, dann ziehen).Vorwärtsrichtung $\pe(A\mid B_k)$ ist gegeben,
    gefragt ist die Rückwärtsrichtung $\pe(B_k\mid A)$ $\Rightarrow$ Bayes.
\end{reg}

\paragraph{Nützliche Regeln (1.6.5).}
$\pe(A^c\mid D)=1-\pe(A\mid D)$;\quad
$\pe(A\cup B\mid D)=\pe(A\mid D)+\pe(B\mid D)-\pe(A\cap B\mid D)$.
\warnung{$\pe(A\cap B\mid D)$ lässt sich \emph{nicht} allein aus $\pe(A\mid D)$
    und $\pe(B\mid D)$ berechnen -- dafür braucht es Zusatzinformation (z.B.\ bedingte
    Unabhängigkeit gegeben $D$).Ferner ist $\pe(C\mid A\cup B)\neq\pe(C\mid A)+\pe(C\mid B)$
    im Allgemeinen.``Ohne Zurücklegen'' $\Rightarrow$ Ziehungen sind abhängig,
    $\pe(A\mid B_k)$ hängt vom ersten Zug ab.}

% ==================================================================

\subsection{Erwartungswert und Varianz}\label{subsec:erwartungswert-und-varianz}
% ==================================================================
Aufg.\ 2.1.4--2.1.6.
\[
    \E{X}=\sum_x x\,\pe(X=x)\quad\text{bzw.}\quad \int x f(x)dx,
\]
\[
    \Var(X)=\E{X^2}-\E{X}^2\geq 0.
\]

\paragraph{Rechenregeln (unverzichtbar!).}
\[
    \E{aX+b}=a\,\E{X}+b,\qquad
    \Var(aX+b)=a^2\,\Var(X).
\]
\[
    \E{X+Y}=\E{X}+\E{Y}\ \text{(immer)},\qquad
    \Var(X+Y)=\Var(X)+\Var(Y)\ \text{(nur falls unkorreliert/unabh.)}.
\]
\warnung{Das $b$ verschwindet in der Varianz, das $a$ wird quadriert
    (z.B.\ $Y=-2X+4$: $\Var(Y)=4\Var(X)$, $\E{Y}=-2\E{X}+4$).Existenz vor der
    Berechnung begründen: bei \emph{endlichem} Bildbereich existieren $\E{}$ und
    $\Var$ stets.Bei $Z=2X+Y$ für $\E{}$ die Linearität nutzen (keine Unabh.\
    nötig); für $\Var(Z)$ dagegen Kovarianz beachten.}

% ==================================================================

\subsection{Die wichtigen Verteilungen im Überblick}\label{subsec:die-wichtigen-verteilungen-im-uberblick}
% ==================================================================

\begin{mytable}
    \begin{tabular}{@{}p{2.8cm} p{4.3cm} p{1.9cm} p{2.0cm} p{2.2cm}@{}}
        \hline
        \textbf{Verteilung} & \textbf{Zähldichte / Dichte} & \textbf{$\E{X}$} &
        \textbf{$\Var(X)$}
        & \textbf{Einsatz}
        \\
        \hline
        Bernoulli $(p)$      & $\pe(1)=p,\ \pe(0)=1-p$  & $p$    & $p(1-p)$ & einzelner Versuch \\
        Binomial $(n,p)$ & $\binom{n}{k}p^k(1-p)^{n-k}$ & $np$ & $np(1-p)$ &
        \#Erfolge, unabh.\ wdh.\\
        Geometrisch $(p)$ & $(1-p)^{k-1}p,\ k\geq 1$ & $\frac1p$ & $\frac{1-p}{p^2}$ &
        Wartezeit bis 1.\ Erfolg \\
        Poisson $(\lambda)$ & $e^{-\lambda}\frac{\lambda^k}{k!}$ & $\lambda$ &
        $\lambda$ & seltene Ereignisse \\
        Gleichvert.\ diskret & $\frac1n$ auf $n$ Werten & Mittel & ---      & Würfel, Laplace   \\
        Gleichvert.\ $[a,b]$ & $\frac{1}{b-a}\mathbf{1}_{[a,b]}$ & $\frac{a+b}2$ &
        $\frac{(b-a)^2}{12}$
        & ``rein zufällig''
        \\
        Normal $(\mu,\sigma^2)$ & $\frac{1}{\sqrt{2\pi}\sigma}e^{-(x-\mu)^2/2\sigma^2}$ &
        $\mu$
        & $\sigma^2$
        & Approximation (GWS)
        \\
        \hline
    \end{tabular}
\end{mytable}
\warnung{Geometrische Verteilung existiert in zwei Konventionen (Zählung ab $1$
    mit $\E{}=\frac1p$ oder ab $0$ mit $\E{}=\frac{1-p}p$).In Aufg.\ 2.2.6 ist die
    ``Nummer des Wurfs'' gemeint $\Rightarrow$ Zählung ab $1$, also $\E{\tau}=\frac1p$
    mit $p=\pe(\text{Augensumme}=5)$.}

% ==================================================================

\subsection{Grenzwertsätze und Approximationen}\label{subsec:grenzwertsatze-und-approximationen}
% ==================================================================
Der Kern von 2.2 und 2.3: Entscheide, \emph{welche} Approximation passt.

\begin{reg}
    \textbf{Faustregel Approximationswahl.}
    \begin{itemize}
        \item $n$ groß,\ $p$ \emph{klein},\ $np=\lambda$ moderat $\Rightarrow$
        \textbf{Poisson} (seltene Ereignisse): Aufg.\ 2.3.4, 2.3.5.
        \item $n$ groß,\ $p$ nicht extrem ($np(1-p)$ groß) $\Rightarrow$
        \textbf{Normal / de Moivre--Laplace}: Aufg.\ 2.2.4, 2.3.6.
        \item Nur Schranke ohne Verteilungsdetails gesucht $\Rightarrow$
        \textbf{Tschebyscheff}: Aufg.\ 2.2.5, 2.2.6.
    \end{itemize}
\end{reg}

\subsubsection{Poisson-Approximation}
\[
    \text{Bin}(n,p)\approx \text{Poi}(\lambda),\quad \lambda=np,\qquad
    \pe(X=k)\approx e^{-\lambda}\frac{\lambda^k}{k!}.
\]

\subsubsection{Satz von de Moivre--Laplace (Normalapproximation)}
Für $S_n\sim\text{Bin}(n,p)$ mit $\mu=np,\ \sigma=\sqrt{np(1-p)}$:
\[
    \pe(a\leq S_n\leq b)\approx
    \Phi\!\Bigl(\frac{b+\frac12-\mu}{\sigma}\Bigr)-
    \Phi\!\Bigl(\frac{a-\frac12-\mu}{\sigma}\Bigr).
\]
\warnung{\textbf{Stetigkeitskorrektur} $\pm\frac12$ nicht vergessen (diskret
    $\to$ stetig).Bei ``$\leq$'' die $+\frac12$, bei ``$\geq$'' die $-\frac12$
    verwenden.Bei \textbf{Überbuchung} (2.2.4) wird die maximale Ticketzahl
    gesucht: Ungleichung $\pe(\text{Erscheinende}>526)\leq 0{,}05$ nach $n$ auflösen,
    $\Phi^{-1}(0{,}95)\approx 1{,}645$ verwenden, Ergebnis \emph{abrunden}.}

\subsubsection{Tschebyscheffsche Ungleichung}
\[
    \pe\bigl(|X-\E{X}|\geq \varepsilon\bigr)\leq \frac{\Var(X)}{\varepsilon^2},
    \qquad
    \pe\bigl(X\geq k\,\E{X}\bigr)\leq \frac{\Var(X)}{(k-1)^2\E{X}^2}.
\]
Damit z.\,B.\ $b_n\to 0$ (schwaches Gesetz großer Zahlen) und Monotonie
$b_{200}<b_{20}$ zeigen.
\warnung{Tschebyscheff liefert nur \emph{obere Schranken}, meist grob.In 2.2.6(g)
    weicht der exakte Wert deutlich von der Abschätzung ab -- das ist erwartet, kein
    Rechenfehler.}

% ==================================================================

\subsection{Statistik: Schätzer}\label{subsec:statistik:-schatzer}
% ==================================================================
Aufg.\ 2.4.4.

\paragraph{Erwartungstreue (unbiased).}
$t$ ist erwartungstreu für $\theta$, falls $\ET[t(X)]=\theta$ für alle
$\theta$.Für das Stichprobenmittel $\bar X=\frac1n\sum X_i$ gilt
$\E{\bar X}=\E{X_1}=\mu$.

\paragraph{Risiko (mittlerer quadratischer Fehler, MSE).}
\[
    R(t)=\E{(t(X)-\theta)^2}
    =\underbrace{\Var(t)}_{\text{Streuung}}+\underbrace{\text{Bias}(t)^2}_{=0\text{ falls erwartungstreu}}.
\]
Für $\bar X$: $R(\bar X)=\Var(\bar X)=\frac{\Var(X_1)}{n}\xrightarrow{n\to\infty}0$
(Konsistenz).
\warnung{Hier \emph{keine} Bernoulli-Annahme -- allgemeine i.\,i.\,d.\ Variablen.
    $\Var(\bar X)=\frac{\sigma^2}{n}$ nutzt Unabhängigkeit (Varianz der Summe =
    Summe der Varianzen).}

% ==================================================================

\subsection{Deskriptive Statistik}\label{subsec:deskriptive-statistik}
% ==================================================================
Aufg.\ 2.4.5.Für Daten $x_1,\dots,x_n$:
\begin{itemize}
    \item \textbf{Arithm.\ Mittel} $\bar x=\frac1n\sum x_i$.
    \item \textbf{Median}: mittlerer Wert der \emph{sortierten} Daten
    (bei geradem $n$: Mittel der beiden mittleren).
    \item \textbf{Modus}: häufigster Wert.
    \item \textbf{Spannweite} $=\max-\min$.
    \item \textbf{Mittlere abs.\ Abw.\ vom Median} $=\frac1n\sum|x_i-\tilde x|$.
    \item \textbf{Varianz} $s^2=\frac1n\sum(x_i-\bar x)^2$, \textbf{SD} $s=\sqrt{s^2}$.
\end{itemize}

\paragraph{Verschiebung um Konstante $c$ (Teil (g)--(i)).}
\[
    \overline{x+c}=\bar x+c,\quad
    \widetilde{x+c}=\tilde x+c,\quad
    s^2_{x+c}=s^2_x\ \text{(unverändert)}.
\]
\warnung{Additive Konstanten verschieben Lagemaße (Mittel, Median), lassen
Streuungsmaße (Varianz, SD, Spannweite) aber \emph{unverändert}.
Daten vor dem Median-Ablesen unbedingt sortieren.Achte auf ein $\frac1n$
    vs.\ $\frac1{n-1}$ je nach Vorlesungskonvention.}

% ==================================================================

\subsection{Konfidenzintervalle}\label{subsec:konfidenzintervalle}
% ==================================================================
Aufg.\ 2.5.4, 2.5.5.Anteilsschätzung $\hat p=\frac{S_n}{n}$, Niveau $1-\alpha$:
\[
    J_n=\Bigl[\hat p - z_{1-\alpha/2}\sqrt{\frac{\hat p(1-\hat p)}{n}},\;
    \hat p + z_{1-\alpha/2}\sqrt{\frac{\hat p(1-\hat p)}{n}}\Bigr].
\]
$z_{1-\alpha/2}=\Phi^{-1}(1-\alpha/2)$ (z.\,B.\ $z_{0{,}975}\approx1{,}96$).
\begin{reg}
    \textbf{Erkennen:} ``Konfidenzintervall zum Niveau $1-\alpha$
    für die Wahrscheinlichkeit/den Anteil'' $\Rightarrow$ obige Formel mit
    $\hat p=\frac{45}{200}=0{,}225$.
\end{reg}
\warnung{Das Niveau bestimmt das Quantil, nicht die Intervalllänge linear:
Verdopplung von $1-\alpha$ verdoppelt \emph{nicht} die Länge (2.5.5(b) ist
\textbf{falsch}).Die Aussage $\pe_p(p\in J_n)\geq 1-\alpha$ (Überdeckung),
    \emph{nicht} $\geq\alpha$ -- Ungleichungsrichtung genau lesen.$p$ ist fest,
    das \emph{Intervall} ist zufällig.}

% ==================================================================

\subsection{Hypothesentests}\label{subsec:hypothesentests}
% ==================================================================
Aufg.\ 2.6.4, 2.6.5.

\paragraph{Ablauf (richtige Reihenfolge, 2.6.5(d)).}
(B) Hypothese \& Alternative $\to$ (E) Teststatistik $\to$ (A) Signifikanzniveau
$\to$ (C) Entscheidungsregel/Verwerfungsbereich $\to$ (D) Stichprobe auswerten
$\to$ (F) Schlussfolgerung.

\paragraph{Einseitiger Anteilstest (2.6.4).}
$H_0: p=p_0\ (=0{,}30)$ gegen $H_1: p<p_0$. Teststatistik unter $H_0$:
\[
    Z=\frac{\hat p-p_0}{\sqrt{p_0(1-p_0)/n}}
    =\frac{S_n-np_0}{\sqrt{np_0(1-p_0)}}\ \overset{H_0}{\approx}\ \mathcal N(0,1).
\]
Verwerfe $H_0$, falls $Z<-z_{1-\alpha}$ (bei $\alpha=0{,}05$: $-1{,}645$).

\paragraph{Fehlerarten \& $p$-Wert (2.6.5).}
\begin{itemize}
    \item \textbf{Fehler 1.\ Art:} $H_0$ verwerfen, obwohl $H_0$ wahr
    (Wahrscheinlichkeit $\leq\alpha$).
    \item \textbf{Fehler 2.\ Art:} $H_0$ beibehalten, obwohl $H_0$ falsch.
    \item \textbf{$p$-Wert:} kleinstes Signifikanzniveau, bei dem $H_0$ gerade noch
    verworfen wird.Wird $H_0$ zu $\alpha=0{,}6$ verworfen, so folgt daraus
    \emph{nicht}, dass sie zu $\alpha=0{,}5$ verworfen wird (kleineres $\alpha$
    = strenger).
\end{itemize}
\warnung{Ein- vs.\ zweiseitig genau bestimmen: ``Wert ist gesunken'' $\Rightarrow$
    \textbf{einseitig} ($H_1: p<p_0$), kritischer Wert $-z_{1-\alpha}$ (nicht
    $-z_{1-\alpha/2}$).Richtung des Verwerfungsbereichs an der Alternative
    ausrichten.``$H_0$ nicht verwerfen'' heißt \emph{nicht} ``$H_0$ bewiesen''.}

\begin{mytable}
    \itshape
    Alle Zahlenwerte (Quantile, $\Phi$-Werte) gemäß\ Vorlesungstabelle einsetzen;
    Rundungs- und Notationskonventionen nach Skript prüfen.
\end{mytable}

